% Simplified Lab Report Format for Chemistry Labs
% Use for: Labs with Purpose [+2], Data Tables, Q&A format, Chemical Structures, ~25 points

\documentclass[
	11pt,
	letterpaper,
]{article}
\usepackage[
	top=1in,
	bottom=1in,
	left=1in,
	right=1in,
]{geometry}

\usepackage[utf8]{inputenc}
\usepackage[T1]{fontenc}
\usepackage{mathpazo}
\usepackage{microtype}
\usepackage{graphicx}
\usepackage{booktabs}
\usepackage{amsmath,amssymb}
\usepackage{float}

% Chemistry packages
\usepackage[version=4]{mhchem}
\usepackage{siunitx}
\sisetup{
	separate-uncertainty=true,
	range-phrase=--,
	range-units=single,
	per-mode=fraction,
	fraction-function=\frac
}
\usepackage{chemfig}

% Formatting
\usepackage{xcolor}
\usepackage{hyperref}
\hypersetup{
	colorlinks=true,
	linkcolor=black,
	urlcolor=blue,
}

\usepackage{titlesec}
\titleformat{\section}{\large\bfseries\sffamily}{\thesection.}{0.5em}{}
\titleformat{\subsection}{\normalsize\bfseries\sffamily}{\thesubsection.}{0.5em}{}

\usepackage{fancyhdr}
\pagestyle{fancy}
\fancyhf{}
\renewcommand{\headrulewidth}{0.4pt}
\fancyhead[L]{\small\textit{Experiment Title}}
\fancyhead[R]{\small\textit{Your Name}}
\fancyfoot[C]{\thepage}

\usepackage{caption}
\captionsetup{
	font=small,
	labelfont=bf,
	justification=justified,
	skip=6pt,
}

% ============================================
% STUDENT INFORMATION - EDIT THESE VALUES
% ============================================
\newcommand{\reporttitle}{Experiment Title}
\newcommand{\reportauthor}{Your Name}
\newcommand{\reportpartner}{Lab Partner Name}
\newcommand{\reportinstructor}{Instructor Name}
\newcommand{\reportcourse}{Course Name}
\newcommand{\reportsection}{Section}
\newcommand{\reportdate}{\today}

\begin{document}

\begin{center}
{\Large\bfseries \reporttitle \par}
\vspace{0.3cm}
{\normalsize \reportauthor \quad \reportpartner \par}
\vspace{0.2cm}
{\small \reportcourse\ (Section \reportsection) \par}
{\small Instructor: \reportinstructor \quad \reportdate \par}
\end{center}

\vspace{0.5cm}

\section{Purpose}

% PURPOSE GOES HERE
% In your own words, describe:
% - The goal of the lab
% - Brief background information to help the reader understand
% - What you are trying to accomplish

\vspace{1em}

% Example: "The purpose of this experiment is to..."

\section{Data Tables}

% DATA TABLES GO HERE
% Include all titled tables with measurements
% Use \begin{table} environment with booktabs

\begin{table}[htbp]
	\centering
	\caption{Descriptive title for this data table.}
	\label{tab:data1}
	\begin{tabular}{@{}lr@{}}
		\toprule
		\textbf{Measurement} & \textbf{Value} \\
		\midrule
		Measurement 1 & \SI{value}{unit} \\
		Measurement 2 & \SI{value}{unit} \\
		Measurement 3 & \SI{value}{unit} \\
		\bottomrule
	\end{tabular}
\end{table}

% Add more tables as needed

\section{Calculations}

% CALCULATIONS GO HERE
% Show all calculations step-by-step
% Use proper equations with \begin{equation} or \begin{align*}

\subsection{Percent Yield Calculation}

\begin{equation}
	\%\ \text{Yield} = \frac{m_{\text{actual}}}{m_{\text{theoretical}}} \times 100\%
\end{equation}

\[
\text{Percent Yield} = \frac{\textit{actual value} \text{ g}}{\textit{theoretical value} \text{ g}} \times 100\% = \textit{result}\%
\]

% Add more calculation subsections as needed

\section{Answers to Questions}

% ANSWERS TO QUESTIONS FROM HANDOUT
% Use \subsection for each question

\subsection{Question 1}

% Answer Question 1 from the handout here

\subsection{Question 2}

% Answer Question 2 from the handout here

\subsection{Question 3}

% Answer Question 3 from the handout here

% Add more questions as needed

\section{Chemical Structures}

% CHEMICAL STRUCTURES GO HERE
% Use \chemfig for structures or include images

\begin{figure}[htbp]
	\centering
	\chemfig{*6(-=-(-C(=[2]O)-[0]OH)=-=)}
	\caption{Structure description.}
	\label{fig:structure1}
\end{figure}

% For complex structures, you can use external images:
% \begin{figure}[htbp]
%   \centering
%   \includegraphics[width=0.5\textwidth]{filename.png}
%   \caption{Structure description.}
% \end{figure}

\subsection{IUPAC Name}

% IUPAC name goes here

\subsection{Functional Groups}

% List and label all functional groups present
% Example:
% \begin{itemize}
%   \item Carboxylic acid group (-COOH)
%   \item Aromatic ring
% \end{itemize}

\section{Conclusion}

% CONCLUSION GOES HERE
% Provide:
% - Brief summary of the lab
% - Main principles involved
% - Key results (percent yield, success assessment)
% - Skills you learned in this lab

\vspace{1em}

% Example structure:
% "In this lab, we investigated [principle]. The key results were [results]. The experiment was [successful/partially successful] as evidenced by [evidence]. Through this lab, I learned [skills]."

\end{document}
